\chapter{LowMC vs AES Analysis}

\section{Why LowMC Should Be Faster (Theory)}

LowMC is widely discussed as a ZK-friendly design because it targets low
multiplicative complexity and minimizes nonlinear gates. Under circuit-centric
reasoning, fewer nonlinear components often imply lower proving cost.

This creates a clear theoretical expectation for circuit-oriented
verifiable-computation settings: LowMC should be competitive, and potentially
faster, when nonlinear gate pressure dominates cost
\cite{albrecht2015ciphers,jaques2019grover}.

That expectation is reasonable at the design level, but it does not by itself
predict end-to-end zkVM runtime. This chapter therefore compares theory against
receipt-backed measurements from the implementation used in this project.

\section{Practical Results}

Measured results in this project do not match the circuit-level expectation. In
the current RISC Zero implementation, LowMC is slower than AES on core proving
metrics.

\paragraph{LowMC vs AES proving-time result.}
Source set: main report benchmark campaign.

Benchmark IDs used: \texttt{aes-r0} and \texttt{lowmc-r0-optimised}.

Trial policy and status: one trial per target; both targets succeeded
(\texttt{success\_rate = 1.0}).

\begin{figure}[t]
\centering
\includegraphics[width=0.78\linewidth]{Figures/report-benchmarks/lowmc_vs_aes_prove_time.png}
\caption{Median proving-time comparison for AES and LowMC from the report benchmark run. In this run, the optimised LowMC target remains substantially slower than AES.}
\label{fig:lowmc-vs-aes-prove}
\end{figure}

In this run, \texttt{lowmc-r0-optimised} is about \(34.85\times\) slower than
\texttt{aes-r0} on proving time (517.209 s vs 14.843 s, both medians from
Table~\ref{tab:cross-algo-comparison}) and about \(40.75\times\) higher on user
cycles.

\paragraph{Baseline vs optimised LowMC result.}
Source set: baseline comparison snapshot.

Benchmark IDs used: \texttt{lowmc-r0} and
\texttt{lowmc-r0-optimised}.

\begin{table}[t]
\centering
\small
\begin{tabular}{lrrrr}
\hline
Metric & Baseline & Optimised & Delta & Delta \% \\
\hline
Prove seconds (median) & 695.050 & 517.209 & -177.841 & -25.59\% \\
Verify seconds (median) & 0.1330 & 0.1065 & -0.0265 & -19.95\% \\
User cycles (median) & 10388894 & 7919919 & -2468975 & -23.77\% \\
Total cycles (median) & 11534336 & 8454144 & -3080192 & -26.70\% \\
\hline
\end{tabular}
\caption{LowMC baseline versus optimised delta (single-trial snapshot, same run).}
\label{tab:lowmc-baseline-vs-optimised}
\end{table}

The optimisation work produced a large measured improvement in this snapshot,
but did not invert the overall ranking against AES.

\section{Comparison with Related Work}

Prior work motivates LowMC by low multiplicative complexity and strong
performance potential in circuit- and MPC-oriented settings
\cite{albrecht2015ciphers,jaques2019grover}. The measurements in this chapter
do not contradict that design motivation; instead, they show a deployment-level
boundary: those gains do not automatically transfer to VM-trace cost in this
implementation stack.

Similarly, ZK-friendly primitive literature often reports constraint-level
benefits for circuit-native constructions. This dissertation's contribution is to
show that VM-level software costs can dominate and reverse expected rankings.

The comparison is therefore best read as a transferability result. LowMC remains
well motivated when the implementation target is a circuit whose cost is closely
related to nonlinear gates. In this dissertation, however, the implementation
target is compiled Rust guest code inside RISC Zero. The cost surface includes
instruction count, memory movement, guest serialization, and trace expansion.
Those costs are visible to the prover even when they are not captured by the
primitive's circuit-level design goal.

\section{Root Cause Analysis}

The likely causes are implementation and VM effects rather than algorithm intent
alone:

\paragraph{Software overhead.}
Abstract gate counts do not include many software costs that appear in a real
zkVM run, such as memory movement, serialization boundaries, and repeated state
transformations in guest code.

\paragraph{Matrix-path pressure.}
LowMC's linear layers require matrix-heavy operations. In this implementation,
those paths create substantial trace and runtime pressure.

\paragraph{Trace expansion effects.}
Repeated bit-level transforms and state updates inflate execution traces, which
raises proving work even when the design target is ``ZK-friendly'' at circuit
level \cite{risczeroTraceDocs}.

\paragraph{VM-vs-circuit optimization gap.}
This project optimizes Rust guest code running on a zkVM; it does not optimize a
custom arithmetic circuit. Circuit-intuition wins therefore do not automatically
transfer to VM-level proving wins.

Supporting pprof evidence makes this explanation more concrete. The screenshots
in Figure~\ref{fig:lowmc-pprof-top-views} come from the LowMC profiling notes and
are illustrative rather than part of the main benchmark campaign. The earlier
unprecomputed path was dominated by in-guest setup work: the top view shows
\texttt{lowmc\_core::rand\_block} at about \(62.75\%\) flat and \texttt{LowMc::new}
on almost the full cumulative path. After moving parameter precomputation to the
host, setup is no longer dominant; input transfer and deserialization inside the
guest become the top entries, especially \texttt{syscall\_2\_nr},
\texttt{sys\_read\_words}, and \texttt{FdReader::read\_words}. This supports the
interpretation that the optimization removed one bottleneck but exposed another
VM-specific boundary cost.

\begin{figure}[p]
\centering
\includegraphics[width=0.94\linewidth]{../notes/images/lowmc-top.png}

\vspace{0.5em}

\includegraphics[width=0.94\linewidth]{../notes/images/lowmc-opt-top.png}
\caption{pprof top views for LowMC before and after the precompute-oriented
input-contract change. The first profile is setup-heavy, while the optimized
profile is dominated by guest input reads and deserialization.}
\label{fig:lowmc-pprof-top-views}
\end{figure}

The later one-time function breakdown from run \texttt{20260318-170118Z} provides
a function-level view of the same issue. Runtime estimates in
Table~\ref{tab:lowmc-pprof-breakdown} are rough attributions derived from RISC
Zero pprof user-cycle shares by multiplying total prove time by each function's
flat-cycle percentage. They are not native CPU timings and should not be read as
precise per-function wall-clock measurements.

\begin{table}[t]
\centering
\scriptsize
\begin{tabular}{p{0.22\linewidth}p{0.36\linewidth}rr}
\hline
Target & pprof entry & Flat share & Approx. prove attribution \\
\hline
\texttt{lowmc-r0} & \texttt{sys\_read\_words} & 60.54\% & 422.7 s \\
\texttt{lowmc-r0} & \texttt{FdReader::read\_words} & 12.81\% & 89.4 s \\
\texttt{lowmc-r0} & serde tuple deserialization & 9.96\% & 69.5 s \\
\texttt{lowmc-r0} & \texttt{multiply\_block\_matrix} & 5.09\% & 35.5 s \\
\hline
\texttt{lowmc-r0-optimised} & \texttt{sys\_read\_words} & 65.43\% & 350.8 s \\
\texttt{lowmc-r0-optimised} & \texttt{FdReader::read\_words} & 13.84\% & 74.2 s \\
\texttt{lowmc-r0-optimised} & serde tuple deserialization & 10.54\% & 56.5 s \\
\texttt{lowmc-r0-optimised} & \texttt{LowMc::encrypt\_rounds} & 0.70\% & 3.8 s \\
\hline
\end{tabular}
\caption{Function-level pprof breakdown for LowMC targets from run
\texttt{20260318-170118Z}. The dominant cost is guest input/deserialization, not
the LowMC round function itself.}
\label{tab:lowmc-pprof-breakdown}
\end{table}

These points are reflected quantitatively: the \(\sim 25.6\%\) LowMC local
improvement in Table~\ref{tab:lowmc-baseline-vs-optimised} still leaves a
large absolute gap versus AES in Figure~\ref{fig:lowmc-vs-aes-prove}.

This supports the root-cause explanation in two ways. First, optimisation of
LowMC internals gives a real improvement, so the implementation path is part of
the cost. Second, the improvement is not enough to approach AES, so the gap is
not explained by one removable inefficiency alone. The remaining difference is
consistent with repeated matrix and bit-level work accumulating across the full
guest trace.

\section{Optimization Attempts and Evidence Trail}

Optimization work targeted expected hotspots, including substitution and matrix
paths. Implemented changes included SWAR-style nonlinear-layer restructuring,
matrix hot-path cleanup, a precompute-oriented host/guest input contract, and
round-count control for controlled experiments.

A concrete receipt-backed evidence snapshot is available in:

\begin{itemize}
\item per-target raw trial JSON for \texttt{lowmc-r0-optimised};
\item corresponding stdout/stderr logs for the same target;
\item aggregate campaign metrics including \texttt{lowmc-r0-optimised};
\item pprof breakdown files from \texttt{20260318-170118Z}, including
\texttt{lowmc-r0.top.flat.txt} and
\texttt{lowmc-r0-optmisied.top.flat.txt} (legacy artifact spelling);
\item pprof screenshots in \texttt{notes/images/lowmc-top.png} and
\texttt{notes/images/lowmc-opt-top.png}.
\end{itemize}

The key lesson is methodological: zkVM optimization must be guided by measured
trace behavior and system overhead, not only by primitive-level design goals.

\section{Takeaway}

The practical conclusion is that LowMC should not be selected for the final
authenticated pipeline in this repository. It is still valuable as a negative
result because it demonstrates exactly why the project benchmarked primitives
inside the target zkVM instead of importing a circuit-level ranking. AES becomes
the better engineering choice not because it is theoretically more ZK-friendly,
but because it is measurably more suitable for this implemented VM workload.
